\begin{table}[!htbp]
\centering
\begin{threeparttable}
\caption{Gaussian versus split-conformal 80\% intervals for AutoARIMA and AutoETS.}
\label{tab:si-conformal}
\small\setlength{\tabcolsep}{4pt}
\begin{tabular}{lccccccccc}
\toprule
Interval & D1 & D2 & D3 & D4 & D5 & D6 & D7 & D8 & D9 \\
\midrule
AutoARIMA, Gaussian & 0.788 & 0.780 & 0.771 & 0.791 & 0.751 & 0.977 & 0.844 & 0.840 & 0.783 \\
AutoARIMA, conformal & 0.613 & 0.591 & 0.614 & 0.605 & 0.601 & 0.654 & 0.651 & 0.614 & 0.602 \\
AutoETS, Gaussian & 0.870 & 0.890 & 0.658 & 0.671 & 0.693 & 0.979 & 0.826 & 0.842 & 0.863 \\
AutoETS, conformal & 0.597 & 0.589 & 0.616 & 0.525 & 0.595 & 0.605 & 0.662 & 0.616 & 0.586 \\
\bottomrule
\end{tabular}
\begin{tablenotes}[flushleft]\footnotesize
\item \textit{Note:} Empirical coverage of the nominal 80\% interval at $n \in \{96, 200\}$. Conformal: split-conformal intervals with two 12-step calibration windows (statsforecast \texttt{ConformalIntervals}); Gaussian: the models' analytic intervals.
\end{tablenotes}
\end{threeparttable}
\end{table}
